% =========================================================
% PART 1 — FOUNDATIONS AND VARIANCE REGIME (FINAL, RECONCILED)
% =========================================================

\section{Introduction and Overview}

The Riemann Hypothesis (RH)~\cite{bombieri} asserts that every nontrivial zero
$\rho = \beta + i\gamma$ of $\zeta(s)$ satisfies $\beta = \tfrac12$.
Equivalently, the analytic “energy’’ of $\zeta(s)$ is balanced across the
critical line.  This paper develops a quantitative, variational formulation
of that balance via a new functional—the
\emph{Spline–Penalized Tail Bound} (SPTB).  

For smoothing width $\Delta>0$ and penalty parameter $\lambda>0$, we measure
the residual between the smoothed tail $H_\sigma(t)$ and its blockwise affine spline
approximation $S_j$ by
\begin{equation}
F_\lambda(H_\sigma;T,\Delta)
  = \sum_{j}\!\int_{I_j}
       \bigl|H_\sigma(t)-S_j(t)\bigr|^2
       + \lambda\,\bigl|\partial_t\!\bigl(H_\sigma(t)-S_j(t)\bigr)\bigr|^2\,dt,
\label{eq:SPTB}
\end{equation}
where $[0,T] = \bigcup I_j$ is partitioned into consecutive blocks of length $\Delta$ and
$S_j$ is the $L^2(I_j)$-best affine fit to $H_\sigma$.

The main analytic results are:
\begin{itemize}
\item[(i)] (\emph{Variance regime}.) If all zeros satisfy $\beta\le\sigma$, then
$F_\lambda(H_\sigma;T,\Delta)=O\!\bigl(T\log T\log\log T\bigr)$ as $T\to\infty$.
\item[(ii)] (\emph{Bias regime / detection}.) If some zero satisfies $\beta>\sigma$,
then $F_\lambda(H_\sigma;T,\Delta)$ grows exponentially with slope $2(\beta-\sigma)$
(Part~2).
\end{itemize}
Together, (i)–(ii) define a finite-window \emph{detector} for violations of RH.

\medskip
\noindent\textbf{Non-equivalence disclaimer.}
We use “detector’’ in the precise sense:
(a) off-line zero $\Rightarrow$ exponential growth (proved);
(b) all zeros on/left of $\sigma$ $\Rightarrow$ polynomial growth (proved);
(c) polynomial boundedness $\Rightarrow$ all zeros on/left of $\sigma$
(\emph{unproved}, the Horocycle Conjecture).  No equivalence with RH is claimed.

\medskip
\noindent\textbf{Abstract (concise).}
\emph{We prove that any zero with $\beta>\sigma$ enforces exponential growth of
$F_\lambda$ as in~\eqref{eq:SPTB}, while if all zeros lie on/left of $\sigma$
the growth is polynomial.  We conjecture (but do not prove) the converse.
Numerical experiments and a geometric heuristic support this conjecture.}

% ---------------------------------------------------------
\section{Scope and Assumptions}

Our analysis rests on four standard inputs for the $L$-function in question
(here stated for $\zeta(s)$; the same template applies \emph{mutatis mutandis}):

\begin{enumerate}
\item[(A1)] \emph{Meromorphic continuation} and the \emph{functional equation.}
\item[(A2)] \emph{Zero counting:}
$N(T)=\tfrac{T}{2\pi}\log\tfrac{T}{2\pi}-\tfrac{T}{2\pi}+O(\log T)$.
\item[(A3)] \emph{Spectral summability:}
$\sum_{|\gamma|\le X}|\rho|^{-2\alpha}=O_\alpha(1)$ for $\alpha>1/2$,
ensuring that $H_\sigma$ as defined in~\eqref{eq:Hsigma} converges in
$L^2_{\mathrm{loc}}$ and its truncations have controlled $L^2$ norms on short intervals.
\item[(A4)] \emph{Montgomery--Vaughan short-interval inequality:} for any
Dirichlet-type sum $D(t)=\sum_{|\gamma|\le X}A_\gamma e^{i\gamma t}$
with distinct real frequencies $\gamma$ of minimal spacing $\delta>0$,
\[
\int_u^{u+\Delta}\!|D(t)|^2\,dt
\;\le\; (\Delta+\delta^{-1})\sum_{|\gamma|\le X}|A_\gamma|^2 .
\]
In every application below the sum is truncated at height
$X=\Delta^{-1}\asymp\log T$ (Theorem~\ref{thm:variance}), where the zero ordinates
have spacing $\delta\gg 1/\log X\asymp 1/\log\log T$; hence
$\delta^{-1}\ll\log X\asymp\log\log T$, which is precisely the additive term recorded
in Lemma~\ref{lem:MV-sieve}.  (The looser bound $\delta^{-1}\ll\log T$ valid for the
full height-$T$ ordinate set is never needed.)
This applies to both $H_\sigma$ and $\partial_t H_\sigma$
(see~\cite{montgomery-vaughan} and Lemma~\ref{lem:MV-sieve}).
\end{enumerate}
No use is made of delicate Euler-product cancellations.  All constants below are
\emph{computable} from the data of (A1)–(A4).

\begin{remark}[Applicability beyond $\zeta$]
Parts~1–2 extend verbatim to Dirichlet $L(s,\chi)$ (primitive $\chi$) and to
automorphic $L$-functions of degree~$d$~\cite{selberg} satisfying analogues of (A1)–(A4).
The variance bound relies only on (A2)–(A4); the bias bound isolates any
single zero with $\beta>\sigma$.  Proofs are stated for $\zeta(s)$ for clarity;
Part~3 validates numerically for $\zeta(s)$, and Appendix~A lists constants.
\end{remark}

% ---------------------------------------------------------
\section{The Harmonic Field \texorpdfstring{$H_\sigma$}{H sigma}}\label{sec:harmonic-field}

Let $\rho=\beta+i\gamma$ run over the nontrivial zeros of $\zeta(s)$
(or more generally of an automorphic $L$-function satisfying (A1)--(A4)).
Fix $\sigma\in[1/2,1)$ and a \emph{smoothing exponent} $\alpha\ge 1$.

\begin{definition}[Harmonic field]\label{def:harmonic-field}
The \emph{harmonic field along $\Re s=\sigma$} is
\begin{equation}
H_\sigma(t)
  \;=\;\sum_\rho \frac{e^{(\beta-\sigma)t}}{|\rho|^\alpha}\,\cos(\gamma t),
  \qquad t\in[0,T].
\label{eq:Hsigma}
\end{equation}
\end{definition}

The weight $|\rho|^{-\alpha}$ serves as a spectral damping factor ensuring absolute
convergence.
By the Riemann--von~Mangoldt formula $N(T)\ll T\log T$,
partial summation gives $\sum_{|\gamma|\le X}|\rho|^{-2\alpha}=O_\alpha(1)$
for every $\alpha>1/2$ (Proposition~\ref{prop:summability} below),
so $H_\sigma(t)$ converges in $L^2_{\mathrm{loc}}$.

\paragraph{Motivation for $\alpha$.}
When $\alpha=1$, the sum $H_\sigma(t)$ closely models the oscillatory part of the
logarithmic derivative $-\zeta'/\zeta(\sigma+it)$ after truncation and smoothing,
with $|\rho|^{-1}$ matching the residue weight of each zero.
Larger $\alpha$ provides faster spectral decay (better convergence) at the cost of
emphasizing low-lying zeros; for the main results we require $\alpha\ge 1$.

\paragraph{Connection to $\zeta(s)$ (heuristic).}
The classical explicit formula~\cite{titchmarsh} gives, for $\sigma>1/2$,
\[
-\frac{\zeta'}{\zeta}(\sigma+it)
=\sum_\rho \frac{1}{(\sigma+it)-\rho} + \text{(polar and $\Gamma$-terms)}.
\]
Each summand satisfies
$\operatorname{Re}\frac{1}{(\sigma+it)-\rho}
=\frac{\sigma-\beta}{(\sigma-\beta)^2+(\gamma-t)^2}$,
which is \emph{not} proportional to $e^{(\beta-\sigma)t}\cos(\gamma t)$;
however, for $|t-\gamma|\ll 1$ and $|\sigma-\beta|$ small, both expressions
share the same exponential sensitivity to $\beta-\sigma$ (the Lorentzian
kernel localizes near $t\approx\gamma$ while the exponential--cosine sum
captures the aggregate oscillatory content).
After inserting the spectral damping weight $|\rho|^{-(\alpha-1)}$
and absorbing the slowly varying Lorentzian prefactor into the blockwise
affine fits, $H_\sigma$ as defined in~\eqref{eq:Hsigma} captures the
oscillatory zero-dependent content of $\zeta'/\zeta$ along $\Re s=\sigma$.
This connection is \emph{motivational}; Theorems~\ref{thm:variance}--\ref{thm:bias}
are proved directly for $H_\sigma$ as defined and require no properties
of $\zeta'/\zeta$ beyond (A1)--(A4).

% ---------------------------------------------------------
\section{Notation and Canonical Regime}

Partition $[0,T]$ into blocks $I_j=[t_j,t_{j+1}]$ of width~$\Delta$ and choose
penalty $\lambda>0$.  Unless otherwise stated we operate in the canonical regime
\begin{equation}
\frac{\kappa_1}{\log T}\le\Delta\le\frac{\kappa_2}{\log T},
\qquad
c_1(\log T)^{-2}\le\lambda\le c_2(\log T)^{-2},
\label{eq:canon}
\end{equation}
with fixed positive constants $\kappa_1,\kappa_2,c_1,c_2$.
Implicit constants in $\ll,\gg,O(\cdot)$ may depend on
$(\alpha,\sigma,\kappa_1,\kappa_2,c_1,c_2)$ but are uniform over~\eqref{eq:canon}.

% ---------------------------------------------------------
\section{Motivation and Background}

Classical equivalence criteria for RH
(Li~\cite{li}, Lagarias~\cite{lagarias}, Beurling~\cite{beurling})
involve global data and infinitely many zeros.  The SPTB functional is
\emph{finite-window} and \emph{local}: it aggregates blockwise amplitude and slope
residuals.  An off-line zero injects a term
$e^{(\beta-\sigma)t}\cos(\gamma t)$ into $H_\sigma$ which no affine projection can remove,
producing an exponential signature in~\eqref{eq:SPTB}.
Conversely, when all zeros lie on/left of~$\sigma$, short-interval
mean-square bounds keep $F_\lambda$ polynomial.

% ---------------------------------------------------------
\section{Affine Projection and Penalization}

Let $S_j$ be the $L^2(I_j)$-best affine fit to $H_\sigma$.
Writing $R = H_\sigma - S$, the functional~\eqref{eq:SPTB}
balances \emph{amplitude} ($\|R\|_{L^2}^2$) and \emph{slope} ($\|R'\|_{L^2}^2$)
on each block.  Affine fits yield transparent constants and suffice
for the sharp lower bounds in the bias regime.
(A $C^2$ cubic-spline variant~\cite{deboor} gives identical asymptotic orders;
see Appendix~B for the derivative-variance constant.)

% ---------------------------------------------------------
\section{Auxiliary Lemmas}\label{sec:aux-lemmas}

We record the technical estimates used throughout the proofs of
Theorems~\ref{thm:variance} and~\ref{thm:bias}.

\begin{lemma}[Affine Projection Lower Bound]\label{lem:affine-lb}
Let $f_{\eta,\omega}(t) = e^{\eta t}\cos(\omega t)$ on $[0,\Delta]$ with $\eta > 0$.
Then
\[
\inf_{a,b}\int_0^{\Delta} |f_{\eta,\omega}(t) - (a+bt)|^2\,dt
\;\ge\; C \cdot
\begin{cases}
\eta^4 \Delta^5, & |\omega| \le \eta/2, \\
\eta^2 \Delta^3, & \eta/2 < |\omega| \le 1/\Delta, \\
\eta^2 \Delta, & |\omega| > 1/\Delta,
\end{cases}
\]
for an absolute constant $C>0$.
\end{lemma}

\begin{proof}
The Gram matrix of the basis $\{1,t\}$ on $[0,\Delta]$ is
\[
G = \begin{pmatrix} \Delta & \Delta^2/2 \\ \Delta^2/2 & \Delta^3/3 \end{pmatrix},
\qquad
G^{-1} = \frac{12}{\Delta^4}
\begin{pmatrix} \Delta^3/3 & -\Delta^2/2 \\ -\Delta^2/2 & \Delta \end{pmatrix}.
\]
The optimal affine approximation error equals
$\|f\|_{L^2}^2 - v^\top G^{-1} v$,
where $v = \bigl(\int_0^\Delta f,\, \int_0^\Delta t\,f\bigr)^\top$.
We compute
\[
v_1=\int_0^\Delta e^{\eta t}\cos(\omega t)\,dt,
\qquad
v_2=\int_0^\Delta t\,e^{\eta t}\cos(\omega t)\,dt.
\]

\textit{Case 1:} $|\omega|\le\eta/2$.
The cosine satisfies $\cos(\omega t)\ge\cos(\eta\Delta/2)\ge 1/2$ on
$[0,\Delta]$ (for $\eta\Delta\le 1$).
Taylor expansion gives $v_1=\Delta+\eta\Delta^2/2+O(\eta^2\Delta^3)$
and $v_2=\Delta^2/2+O(\eta\Delta^3)$.
The projection captures most of $\|f\|^2=\Delta+\eta\Delta^2+O(\eta^2\Delta^3)$,
leaving a residual $\ge C\eta^4\Delta^5$: the leading non-affine term of
$e^{\eta t}-(\hat a+\hat b t)$ is the \emph{quadratic} $\tfrac12\eta^2 t^2$, whose
$L^2(I)$ projection residual against $\mathrm{span}\{1,t\}$ equals
$\bigl(\tfrac{\eta^2}{2}\bigr)^2\cdot\tfrac{\Delta^5}{180}=\Theta(\eta^4\Delta^5)$.

\textit{Case 2:} $\eta/2<|\omega|\le 1/\Delta$.
Since $|\omega|\Delta\le 1$, the cosine completes at most one oscillation on
$[0,\Delta]$; hence $\|f\|^2=\int_0^\Delta e^{2\eta t}\cos^2(\omega t)\,dt\ge c\,\Delta$.
For the moments, write $e^{\eta t}=1+\eta t+\tfrac12\eta^2t^2+O(\eta^3\Delta^3)$.
Since $\cos(\omega t)$ is bounded and slowly varying ($|\omega|\Delta\le 1$),
\[
v_1=\int_0^\Delta (1+\eta t+\cdots)\cos(\omega t)\,dt = O(\Delta),
\quad
v_2=\int_0^\Delta t(1+\eta t+\cdots)\cos(\omega t)\,dt = O(\Delta^2).
\]
The Gram projection satisfies
$v^\top G^{-1}v
=\frac{12}{\Delta^4}\bigl(\tfrac{\Delta^3}{3}v_1^2
-\Delta^2 v_1 v_2+\Delta\,v_2^2\bigr)
=O(\Delta)$.
The quadratic term $\tfrac12\eta^2 t^2\cos(\omega t)$ is not affine;
its $L^2$ projection residual onto $\{1,t\}$ on $[0,\Delta]$ is
$\Theta(\eta^2\Delta^3)$, giving
$\|f\|^2-v^\top G^{-1}v\ge c'\eta^2\Delta^3$.

\textit{Case 3:} $|\omega|>1/\Delta$.
The integrals $v_1,v_2$ involve a rapidly oscillating factor
$\cos(\omega t)$ with $|\omega|\Delta>1$.
By the Riemann--Lebesgue lemma,
$|v_1|\le 2/|\omega|$ and $|v_2|\le 2\Delta/|\omega|+2/\omega^2$.
Hence
\[
v^\top G^{-1}v
=\frac{12}{\Delta^4}\bigl(\tfrac{\Delta^3}{3}v_1^2
-\Delta^2 v_1 v_2+\Delta\,v_2^2\bigr)
=O\bigl(\omega^{-2}/\Delta\bigr).
\]
For $\|f\|^2$, writing $\cos^2(\omega t)=\tfrac12+\tfrac12\cos(2\omega t)$
gives $\int_0^\Delta e^{2\eta t}\cos^2(\omega t)\,dt
=\tfrac12\int_0^\Delta e^{2\eta t}\,dt+O(1/|\omega|)
\ge c\,\Delta$
(since the oscillatory part averages out).
The non-affine contribution from $e^{\eta t}$ has
$L^2$ energy $\Theta(\eta^2\Delta)$, and the Gram projection captures
at most $O(\omega^{-2}/\Delta)\ll\eta^2\Delta$ for $|\omega|>1/\Delta$,
so the residual is $\ge C\eta^2\Delta$.
\end{proof}

\begin{lemma}[Derivative Penalty]\label{lem:derivative-penalty}
For $f_{\eta,\omega}(t)=e^{\eta t}\cos(\omega t)$ on $I=[a,a+\Delta]$
with $\eta\Delta\le 1/2$ \emph{and} $|\omega|\Delta\ge 1$,
\[
\inf_b \int_I |(f_{\eta,\omega}(t)-bt)'|^2\,dt
\;\ge\; c\, e^{2\eta a}\,\omega^2\,\Delta
\;\ge\; c\, e^{2\eta a}\,\eta^2\min\{\Delta,1/|\omega|\},
\]
for an absolute constant $c>0$.
\end{lemma}

\begin{remark}[Role of the oscillation hypothesis $|\omega|\Delta\ge1$]\label{rmk:osc-hypothesis}
The hypothesis $|\omega|\Delta\ge1$ (at least one radian of phase per block) is
\emph{necessary}.  As $\omega\to0$ the derivative $f'$ becomes nearly affine on $I$, and
the left-hand side degenerates to $\Theta\!\bigl(e^{2\eta a}\eta^4\Delta^3\bigr)$, which is
smaller than $\eta^2\Delta$ by the factor $(\eta\Delta)^2$; no lower bound of the form
$c\,\eta^2\min\{\Delta,1/|\omega|\}$ can then hold with an absolute $c>0$.
In the detection regime of Part~2 the relevant frequency is a zero ordinate
$\omega=\gamma$ with $|\gamma|\Delta\gg1$ throughout the canonical scaling, so the
hypothesis is amply satisfied; the complementary low-frequency case is handled by the
amplitude bound (Lemma~\ref{lem:affine-lb}), which needs no oscillation hypothesis.
\end{remark}

\begin{proof}
The minimising constant is the mean $\bar{f'}=\Delta^{-1}\int_I f'$, so the left-hand side
is the variance of $f'$ on $I$:
\[
\inf_b \int_I |f'-b|^2\,dt
= \int_I |f'|^2 - \Delta\,\bar{f'}^{\,2}.
\]
Since $\eta\Delta\le\tfrac12$ the weight obeys $e^{\eta a}\le e^{\eta t}\le e^{1/2}e^{\eta a}$
on $I$, so it may be frozen at $t=a$ at the cost of an absolute factor; it therefore suffices
to bound the variance of $g(t)=\eta\cos(\omega t)-\omega\sin(\omega t)$, whose amplitude is
$\sqrt{\eta^2+\omega^2}\ge|\omega|$.  A direct evaluation gives
$\int_I g^2=\tfrac12(\eta^2+\omega^2)\Delta+O\!\bigl((\eta^2+\omega^2)/|\omega|\bigr)$
and $\bigl|\int_I g\bigr|=O(1)$.  Because $|\omega|\Delta\ge1$ the main term dominates the
subtracted mean, so the variance of $g$ is $\gg\omega^2\Delta$; restoring the frozen weight
gives $\inf_b\int_I|f'-b|^2\ge c\,e^{2\eta a}\omega^2\Delta$.
Finally $|\omega|\Delta\ge1$ together with $\eta\Delta\le\tfrac12$ forces $|\omega|\ge2\eta$
and $\min\{\Delta,1/|\omega|\}=1/|\omega|$, whence
$\omega^2\Delta=\omega\,(\omega\Delta)\ge\omega\ge\eta^2/|\omega|
=\eta^2\min\{\Delta,1/|\omega|\}$, which is the second inequality.
\end{proof}

\begin{lemma}[Montgomery--Vaughan Large Sieve]\label{lem:MV-sieve}
For $I=[u,u+\Delta]$ and real frequencies $\gamma$ with spacing
$\gg 1/\log T$,
\[
\int_I \Bigl|\sum_{|\gamma|\le X} A_\gamma\, e^{i\gamma t}\Bigr|^2 dt
\;\ll\; (\Delta + \log X)\sum_{|\gamma|\le X} |A_\gamma|^2.
\]
\end{lemma}

\begin{proof}
This is a direct specialization of the Montgomery--Vaughan large sieve
inequality~\cite{montgomery-vaughan}, using standard zero spacing estimates
$\gg 1/\log T$; see also Appendix~\ref{app:dirichlet} for details on
sieve applicability.
\end{proof}

\begin{proposition}[Summability]\label{prop:summability}
For $\alpha > 1/2$,
\[
\sum_{|\gamma|\le X} |\rho|^{-2\alpha} = O_{\alpha}(1),
\]
uniformly in $X$.
\end{proposition}

\begin{proof}
Partial summation against $dN(T)$, with
$N(T)=\frac{T}{2\pi}\log\frac{T}{2\pi} - \frac{T}{2\pi}+O(\log T)$,
shows convergence of $\sum |\gamma|^{-2\alpha}$ for $2\alpha>1$.
\end{proof}

\begin{lemma}[High-Frequency Tail]\label{lem:high-freq}
Let $X=\Delta^{-1}=\log T/(2\pi)$. For $\alpha\ge 1$ and $\beta\le\sigma$,
\[
\int_I \Bigl|\sum_{|\gamma|>X}
\frac{e^{(\beta-\sigma)t}\cos(\gamma t)}{|\rho|^\alpha}\Bigr|^2 dt
\;\ll\; \Delta \cdot X^{-(2\alpha-1)}\log X.
\]
\end{lemma}

\begin{proof}
Apply Cauchy--Schwarz: split $\sum |\rho|^{-2\alpha}$ from zero density
estimates.  The partial sum over $|\gamma|>X$ contributes $O(X^{-(2\alpha-1)})$
by Proposition~\ref{prop:summability} and partial summation;
the factor $\log X$ accounts for zero-counting growth.
\end{proof}

% ---------------------------------------------------------
\section{Variance Regime: On-Line Zeros}

When all zeros satisfy $\beta\le\sigma$, the local slope energy of $H_\sigma$
is bounded by Montgomery–Vaughan short-interval inequalities together with
(A2)–(A3):

\begin{theorem}[Variance Regime]\label{thm:variance}
Assume \textup{(A1)–(A4)}.
For $\sigma\ge\tfrac12$, $\lambda\asymp(\log T)^{-2}$, and
$\kappa_1/\log T \le \Delta \le \kappa_2/\log T$, one has
\begin{equation}
F_\lambda(H_\sigma;T,\Delta)
  = O_{\sigma,\alpha}\!\bigl(T\log T\log\log T\bigr).
\label{eq:variance}
\end{equation}
\end{theorem}

\noindent\emph{Unconditionality.}
Theorem~\ref{thm:variance} depends only on (A2)–(A4) and holds independently
of RH.  Constants are uniform over the canonical regime~\eqref{eq:canon}.

\begin{proof}[Proof of Theorem~\ref{thm:variance}]
Write $H_\sigma = H_{\le X} + H_{>X}$ with $X = \Delta^{-1} = \log T/(2\pi)$.
The precise location of the cut is immaterial: any $X$ with
$\Delta^{-1}\le X\le(\log T)^2$ yields the same final bound, the low- and
high-frequency estimates below being valid throughout that range;
Appendix~\ref{app:A} carries out the identical split at $\Gamma_0=(\log T)^2$.
We use the competitor $S_j = 0$ on each block (which only overestimates
the infimum over affine fits).

\medskip\noindent
\textit{Low-frequency variance via large sieve.}
Since $e^{(\beta-\sigma)t}\le 1$ uniformly in $t$ when all $\beta\le\sigma$,
Lemma~\ref{lem:MV-sieve} yields on any block~$I_j$:
\[
\int_{I_j} |H_{\le X}|^2\,dt
\;\ll\; (\Delta + \log X)\sum_{|\gamma|\le X} |\rho|^{-2\alpha}.
\]
By Proposition~\ref{prop:summability}, the sum is $O(1)$ for $\alpha > 1/2$. Hence
\[
\int_{I_j} |H_{\le X}|^2\,dt \;\ll\; \log\log T.
\]

\medskip\noindent
\textit{High-frequency tail.}
By Lemma~\ref{lem:high-freq} on each block~$I_j$,
\[
\int_{I_j} |H_{>X}|^2\,dt
\;\ll\; \Delta \cdot X^{-(2\alpha-1)}\log X
\;\ll\; (\log T)^{-1}\log\log T
\]
for $\alpha\ge 1$.

\medskip\noindent
\textit{Derivative contribution (low-frequency).}
Differentiating $H_{\le X}$:
\[
\partial_t H_{\le X}
= \sum_{|\gamma|\le X}
\frac{e^{(\beta-\sigma)t}}{|\rho|^\alpha}
\bigl((\beta-\sigma)\cos(\gamma t)-\gamma\sin(\gamma t)\bigr).
\]
Applying Lemma~\ref{lem:MV-sieve} to $\partial_t H_{\le X}$,
\[
\int_{I_j} |\partial_t H_{\le X}|^2\,dt
\;\ll\; (\Delta + \log X)
\sum_{|\gamma|\le X}\frac{(\beta-\sigma)^2+\gamma^2}{|\rho|^{2\alpha}}.
\]
When $\alpha=1$ the sum is $\ll N(X)\ll X\log X\ll\log T\log\log T$;
when $\alpha>1$ the sum converges uniformly.
In both cases,
\[
\int_{I_j} |\partial_t H_{\le X}|^2\,dt \;\ll\; \log T(\log\log T)^2.
\]
Multiplying by $\lambda = (\log T)^{-2}$ contributes $\ll (\log\log T)^2/\log T$
per block.

\medskip\noindent
\textit{Derivative contribution (high-frequency).}
For the tail $\partial_t H_{>X}$ the coefficient sum
$\sum_{|\gamma|>X}\gamma^2|\rho|^{-2\alpha}$ diverges when $\alpha=1$, so a coefficient
bound is insufficient and the short-interval cancellation of
Lemma~\ref{lem:MV-sieve} is essential.  Applying it to $\partial_t H_{>X}$ on $I_j$
(equivalently, the per-zero estimate~\eqref{eq:high} of Appendix~\ref{app:A}) gives
\[
\int_{I_j}|\partial_t H_{>X}|^2\,dt
\;\ll\;\sum_{|\gamma|>X}\frac{\gamma^2}{|\rho|^{2\alpha}}
\min\!\Bigl\{\Delta,\frac{1}{\gamma^2\Delta}\Bigr\}
\;=\;\frac{1}{\Delta}\sum_{|\gamma|>X}\frac{1}{|\rho|^{2\alpha}}
\;\ll\;\frac{X^{-(2\alpha-1)}\log X}{\Delta}\;\ll\;\log\log T,
\]
since $|\gamma|>X=\Delta^{-1}$ forces $\min\{\Delta,1/(\gamma^2\Delta)\}=1/(\gamma^2\Delta)$,
while $\sum_{|\gamma|>X}|\rho|^{-2\alpha}\ll X^{-(2\alpha-1)}\log X$ by
Lemma~\ref{lem:high-freq} and $\Delta^{-1}=X$ (so the bound is $X^{2-2\alpha}\log X\ll\log\log T$
for $\alpha\ge1$).  Multiplying by $\lambda=(\log T)^{-2}$ this is
$\ll\log\log T/(\log T)^2$ per block, dominated by the low-frequency derivative term.

\medskip\noindent
\textit{Summation over blocks.}
Combining all terms, each block contributes $\ll \log\log T$ to $F_\lambda$.
There are $J \ll T/\Delta \ll T\log T$ blocks. Hence
\[
F_\lambda(H_\sigma;T,\Delta) \;\ll\; T\log T\log\log T. \qedhere
\]
\end{proof}

% ---------------------------------------------------------
\section{Robustness in the Penalty Parameter \texorpdfstring{$\lambda$}{lambda}}

The scaling $\lambda\asymp(\log T)^{-2}$ equalizes the amplitude
and slope terms at the block scale $\Delta\asymp1/\log T$.
Both the variance bound~\eqref{eq:variance} and the bias lower bounds
(Part~2) remain valid for any
$\lambda\in[c_1(\log T)^{-2},\,c_2(\log T)^{-2}]$,
changing only multiplicative constants.
Numerically, the exponential slope in the bias regime varies by $<0.2\%$
across a $16\times$ range of~$\lambda$, confirming detector stability.

% ---------------------------------------------------------
\section{Heuristic Interpretation}

Formally, $F_\lambda$ behaves like a curvature-regularized Fisher information:
bounded growth corresponds to finite “information curvature,” while an
off-line zero induces a curvature singularity with exponential rate
$2(\beta-\sigma)$—made precise analytically in Part~2.
This view motivates but does not replace the rigorous arguments below.

